Este apêndice lista as verificações que decidem a aprovação de cada experimento. Elas estão escritas como asserções (\texttt{assert}) nas últimas células do \textit{notebook} de experimentos \cite{aeroledger_experimentos}, separadas do código que produz os resultados, e foram registradas no repositório antes da execução. Cada asserção compara um resultado com um valor esperado e, quando é o caso, com a margem definida na Seção~\ref{sec:experimentos}.

\begin{longtable}{l p{0.82\linewidth}}
    \caption{Verificações de cada experimento.}\label{tab:ape-testes}\\
    \toprule
    & Verificação \\
    \midrule
    \endfirsthead
    \toprule
    & Verificação \\
    \midrule
    \endhead
    \bottomrule
    \endfoot
    E1 & Doze campos mudam entre as edições de 2022 e 2026: o comprimento da pista, o comprimento da zona de parada, o comprimento e a largura da faixa e as oito distâncias declaradas. \\
    E1 & O evento que registra a edição de 2026 não traz causa. \\
    \midrule
    E2 & Cada grandeza calculada difere da publicada dentro da sua margem: 0,919~m para o comprimento e as distâncias declaradas, 1,419~m para a ASDA~16, 0,014$^\circ$ para os azimutes e 1,009$^\circ$ para os rumos magnéticos. \\
    E2 & Nenhum valor derivado está gravado no banco. \\
    E2 & A remedição grava exatamente dois eventos, o \texttt{BASELINE} do levantamento e o \texttt{DELTA} do aeródromo, e este altera apenas o identificador do referencial. \\
    E2 & A troca de referencial tem o levantamento como causa. \\
    E2 & Nenhuma \textit{feature} com posição (pista, cabeceiras e zona de parada) é reescrita. \\
    E2 & Os vértices, o comprimento e o azimute calculados pela biblioteca coincidem com o gabarito de Vincenty a menos de 0,5~mm. \\
    \midrule
    E3 & Em $t_a$ vigoram o referencial e o comprimento antigos, e, depois da remedição, os novos. \\
    E3 & A interdição existe nos três instantes dentro do seu intervalo e não existe nos três instantes fora dele. \\
    E3 & O ciclo de vida da interdição é exatamente um \texttt{BASELINE} seguido de um \texttt{VOID}, nos instantes em que foram escritos. \\
    E3 & O envelope da troca de referencial informa o autor, o motivo e a causa, e o levantamento que a causou tem autor. \\
    E3 & As distâncias declaradas recalculadas para $t_a$ são idênticas às calculadas no próprio $t_a$. \\
    \midrule
    E4 & A alteração de um evento é apontada no próprio evento, e a exclusão e a inserção, no evento seguinte. \\
    E4 & A reescrita da cadeia até o fim não é apontada pela verificação interna. \\
    E4 & Os quatro ataques são detectados pela comparação com a âncora. \\
    \midrule
    E5 & Os dois produtos apresentam as mesmas distâncias declaradas em $t_a$ e depois da remedição, e essas distâncias mudam entre os dois instantes. \\
    E5 & O verificador independente declara a cadeia íntegra e chega à mesma cabeça que a biblioteca. \\
    E5 & O comprimento calculado pelo verificador difere do da biblioteca em menos de 0,5~mm. \\
    E5 & No arquivo alterado, o verificador aponta o evento cuja largura foi modificada. \\
    \midrule
    E6 & Em todas as fronteiras de corte e nos três instantes testados, o consumidor apresenta os mesmos estados e as mesmas distâncias que o banco. \\
    E6 & O banco, cortado no meio da carga, informa a ASDA~16 igual à TORA~16, e nenhum corte do consumidor cai dentro da carga. \\
    E6 & O consumidor que assina apenas a cabeceira~16 aplica um fecho menor e chega aos mesmos estados e ao mesmo \textit{hash} dos derivados. \\
    E6 & Antes de receber a restrição nova, o consumidor recusa a leitura; depois, responde incluindo-a, e a resposta não muda com o banco fechado. \\
    E6 & Os seis consumidores chegam ao mesmo \textit{hash} no mesmo limite. \\
    E6 & Os quatro ataques no meio do tópico interrompem o consumo e deixam o consumidor fiel ao banco. \\
    E6 & O corte do fim do tópico não interrompe o consumo, e o evento falso anexado ao fim é aplicado. \\
\end{longtable}
